\section{Group Discussion and Conclusion}

\subsection{Group Discussion}
Sau khi hoàn thành bản tích hợp hiện tại, nhóm thống nhất một số điểm kỹ thuật quan trọng:
\begin{itemize}
    \item Thiết kế theo kiến trúc phân lớp bao gồm Input, FSM, Service và Driver giúp mở rộng tính năng nhanh chóng và giảm rủi ro chồng chéo logic.
    \item Sử dụng semaphore và mutex là yêu cầu bắt buộc trong môi trường đa tác vụ, đặc biệt tại các vùng có sự tranh chấp tài nguyên như GPIO, danh sách relay hoặc việc ghi nhật ký hệ thống.
    \item Việc triển khai TinyML với các mô hình định dạng int8 là lựa chọn phù hợp cho ESP32-S3 vì giúp tối ưu hóa bộ nhớ nhưng vẫn đảm bảo khả năng ra quyết định điều khiển hiệu quả dựa trên môi trường.
    \item AP dashboard đóng vai trò then chốt cho vận hành thực tế, vì cho phép cấu hình tại chỗ khi thiết bị chưa vào mạng STA.
\end{itemize}

\subsection{Conclusion}
Đồ án đã hiện thực hệ thống IoT plant-care trên ESP32-S3 và đáp ứng bộ yêu cầu chính gồm điều khiển LED/NeoPixel/OLED theo ngữ cảnh cảm biến, web server AP điều khiển thiết bị, suy luận TinyML on-device, và publish telemetry lên CoreIOT ở STA mode. Ngoài yêu cầu bắt buộc, nhóm bổ sung OTA firmware update, manual mode qua IR remote, và logging framework thống nhất để tăng khả năng vận hành/debug.

Kết quả thử nghiệm cho thấy luồng chức năng end-to-end hoạt động ổn định trong điều kiện lab. Phần TinyML đã có metric huấn luyện/validation tốt và chạy thành công trên thiết bị thật. Trong các phiên bản tiếp theo, nhóm ưu tiên:
\begin{itemize}
    \item xây dựng bộ đánh giá tự động cho độ chính xác trên phần cứng thông qua việc tính toán ma trận nhầm lẫn trực tiếp trên bảng mạch,
    \item tái cấu trúc mã nguồn nhằm giảm sự phụ thuộc vào các biến trạng thái toàn cục theo hướng kiến trúc dựa trên ngữ cảnh,
    \item mở rộng thêm các thiết bị chấp hành, chẳng hạn như máy bơm thực tế, để tận dụng tối đa không gian hành động của mô hình.
\end{itemize}
